\scanpage{6}
\mspage{3}

\margin{\unsure{*|4}}

Je me dispense d'écrire le carré cocartésien dual, quand $A$ a les $\underset{\unsure{-|\to}}{\lim}$ pertinentes, y compris $\unsure{\varphi|\emptyset}$ objet initial. Il semble qu'elle n'ait pas servi en pratique. Les foncteurs adjoints et \unsure{coadjoints|multiadjoints} \unsure{classiques|d'usage} utiles sont \deleted{\unsure{écrits}} \inserted{indiqués} ci-dessous (où les ``\unsure{peu|pas} utiles'' sont entre crochets)

\[
\begin{aligned}
&[i_{!}]\,i^{*}\,i_{*}\,i^{!}\\
&[j^{?}]\,j_{!}\,j^{*}\,j_{*}\ ,
\end{aligned}
\]

i.e. les foncteurs $i_{!}$ et $j^{?}$ n'ont pas servi jusqu'à \unsure{maint.}

f) Cas des catégories $\Delta_{i}$ (simplexes types, $i\in N$), et des foncteurs entre ceux-ci (applications semi-simpliciales). \deleted{\unsure{Soient}} Il semblerait que pour une telle $f$, \deleted{\unsure{quel que soit $f$}} \inserted{si $A$} \deleted{\unsure{admet un}} est pointée, il y ait des \uline{bi-adjoints} à la fois à g. et à dr. De plus, pour $\Delta_{i}$ et $\Delta_{j}$ fixés,

Quand on regarde toutes les $f:\Delta_{i}\unsure{-|\to}\Delta_{j}$ et toutes les $g:\Delta_{j}-\Delta_{i}$, que les suites \inserted{de toutes} \deleted{adjoints} multiadjoints pour \inserted{chacun} des $f^{*},g^{*}$ correspondants sont toutes les mêmes. Ça a l'air vrai du moins si $j=i+1$ (et je l'ai vérifié pour $i=0$, $i=1$), on doit

trouver une ``suite d'adjonction'' de longueur $2i+7\inserted{{}=(i+1)+(i+2)+2+2}$. De façon \deleted{\unsure{$2i+7$}} précise, la suite comporte \deleted{\unsure{deux}} \inserted{quatre} termes ``extrêmes'' (les deux premiers, les deux derniers) et les $\unsure{(i+1)|(2i+1)}+(i+2)$ termes intermédiaires correspondent aux foncteurs de la forme $f^{*}$ ou $g^{*}$, où $f\inserted{{}:\Delta_{i+1}-\Delta_{i}}$ est une des $(i+1)$ applications de dégénérescence, et $g:\Delta_{i}-\Delta_{i+1}$ une des $i+2$ applications bord.

\figure{\begin{tikzpicture}\draw (0,.16)--(.09,0);\end{tikzpicture}}

% archive sheet number: 5
